\begin{theorem}[一参数子群的频率参数化与实投影读出]\label{thm:cdim-arithmetic-singular-ring-one-parameter-subgroups}
\leanverified{paper\_cdim\_arithmetic\_singular\_ring\_one\_parameter\_subgroups}
设 \(S\) 为有限素数集，\(|S|=s\ge 2\)，并令算术奇异环 \(\mathfrak{R}_S=\Sigma_S\times_{\TT}\TT^S\) 如 \eqref{eq:cdim-arithmetic-singular-ring-definition}。
则存在自然同构
\[
\Hom_{\mathrm{cont}}(\RR,\mathfrak{R}_S)\ \cong\ \RR^{S}.
\]
更具体地，对任意一参数子群 \(\gamma:\RR\to\mathfrak{R}_S\)，其在 \(\TT^S\) 上的投影可唯一写为
\begin{equation}\label{eq:cdim-arithmetic-singular-ring-exp-orbit}
\mathrm{pr}_2(\gamma(t))=\bigl(e^{\mathrm{i}\nu_p t}\bigr)_{p\in S}
\qquad(t\in\RR)
\end{equation}
其中 \(\nu=(\nu_p)_{p\in S}\in\RR^S\) 唯一决定 \(\gamma\)；
并且 \(\Sigma_S\) 坐标满足纤维积兼容条件
\[
\pi_S\bigl(\mathrm{pr}_1(\gamma(t))\bigr)
=
\mu_S\bigl(\mathrm{pr}_2(\gamma(t))\bigr)
=
e^{\mathrm{i}\left(\sum_{p\in S}\nu_p\right)t}.
\]
在 \(S=\{p,q\}\) 的情形，取整数频率对 \((a,b)\in\ZZ^2\) 与相位 \(\delta\in\RR\)，则平移轨道
\[
t\longmapsto \bigl(e^{\mathrm{i}(at+\delta)},\,e^{\mathrm{i}bt}\bigr)\in\TT^2
\]
为一参数子群 \(t\mapsto (e^{\mathrm{i}at},e^{\mathrm{i}bt})\) 的陪集。
选取任意 \(x_{\delta}\in\Sigma_{\{p,q\}}\) 满足 \(\pi_{\{p,q\}}(x_\delta)=e^{\mathrm{i}\delta}\)，并置
\(
g_{\delta}:=(x_{\delta},(e^{\mathrm{i}\delta},1))\in\mathfrak{R}_{\{p,q\}}
\)，
则陪集 \(g_{\delta}\cdot\gamma_{(a,b)}(\RR)\subset\mathfrak{R}_{\{p,q\}}\) 在 \(\TT^2\) 上的投影恰为上式；
对其再取坐标实部得到 Lissajous 轨迹
\[
\bigl(\Re(e^{\mathrm{i}(at+\delta)}),\,\Re(e^{\mathrm{i}bt})\bigr)
=\bigl(\cos(at+\delta),\,\cos(bt)\bigr).
\]
对给定 \(\delta\)，提升元 \(x_{\delta}\) 的选择在 \(\ker(\pi_{\{p,q\}})\) 的平移下不唯一；
该自由度不改变 \(\TT^2\) 轨道，仅改变陪集在 \(\Sigma_{\{p,q\}}\) 方向的“隐蔽环向”位置。
\end{theorem}

\begin{proof}
由定理 \ref{thm:cdim-arithmetic-singular-ring-dual-connected}，\(\Gamma_S:=\widehat{\mathfrak{R}_S}\) 为秩 \(s\) 的无挠离散群，并具有表示
\[
\Gamma_S\cong (\ZZ[S^{-1}]\oplus \ZZ^{S})/\langle(1,-\mathbf{1}_S)\rangle.
\]
Pontryagin 对偶给出自然同构
\[
\Hom_{\mathrm{cont}}(\RR,\mathfrak{R}_S)\ \cong\ \Hom(\Gamma_S,\widehat{\RR}),
\]
而 \(\widehat{\RR}\cong\RR\)（取离散拓扑）。
因此 \(\Hom_{\mathrm{cont}}(\RR,\mathfrak{R}_S)\cong \Hom(\Gamma_S,\RR)\cong \RR^S\)。
\par
为得到 \eqref{eq:cdim-arithmetic-singular-ring-exp-orbit} 的显式形式，令 \((e_p)_{p\in S}\) 表示 \(\ZZ^{S}\) 的标准基。
给定 \(\nu=(\nu_p)_{p\in S}\in\RR^{S}\)，定义同态 \(\ell:\ZZ[S^{-1}]\oplus\ZZ^{S}\to\RR\) 为
\[
\ell(a,(n_p)_{p\in S})
:=a\Bigl(\sum_{p\in S}\nu_p\Bigr)+\sum_{p\in S}n_p\nu_p.
\]
则 \(\ell(1,-\mathbf 1_S)=0\)，故 \(\ell\) 因子化为 \(\Gamma_S\to\RR\)，从而对应一参数子群 \(\gamma_{\nu}:\RR\to\mathfrak{R}_S\)。
在 \(\TT^S\) 坐标上，字符 \(e_p\) 的取值沿 \(\gamma_\nu\) 等于 \(t\mapsto e^{\mathrm{i}\ell(e_p)t}=e^{\mathrm{i}\nu_pt}\)，即得到 \eqref{eq:cdim-arithmetic-singular-ring-exp-orbit}；
而关系元 \((1,-\mathbf{1}_S)=0\) 等价于
\(
\ell(1)=\sum_{p\in S}\ell(e_p)
\)，
从而给出 \(\pi_S\circ \mathrm{pr}_1=\mu_S\circ \mathrm{pr}_2\) 下的全局相位频率 \(\sum_{p\in S}\nu_p\)。
反向唯一性由对偶同态的唯一性与指数轨道确定性直接推出。
\par
在 \(S=\{p,q\}\) 情形，取 \(\nu=(a,b)\) 并记相应一参数子群为 \(\gamma_{(a,b)}\)。
由 \(\pi_{\{p,q\}}\) 的满射性可取 \(x_{\delta}\in\Sigma_{\{p,q\}}\) 使 \(\pi_{\{p,q\}}(x_\delta)=e^{\mathrm{i}\delta}\)，从而 \(g_{\delta}\in\mathfrak{R}_{\{p,q\}}\)。
则 \(\mathrm{pr}_2(g_{\delta}\cdot\gamma_{(a,b)}(t))=(e^{\mathrm{i}\delta},1)\cdot(e^{\mathrm{i}at},e^{\mathrm{i}bt})=(e^{\mathrm{i}(at+\delta)},e^{\mathrm{i}bt})\)，取坐标实部即得 Lissajous 轨迹。
若改取 \(x'_{\delta}=k\cdot x_{\delta}\)（\(k\in\ker(\pi_{\{p,q\}})\)），则所得陪集在 \(\TT^2\) 上的投影不变，而在 \(\Sigma_{\{p,q\}}\) 方向作平移；该非唯一性因此精确由 \(\ker(\pi_{\{p,q\}})\) 参数化。证毕。
\end{proof}

\begin{proposition}[素频率和约束下的特征评价公式]\label{prop:cdim-arithmetic-singular-ring-prime-frequency-character-eval}
沿用定理 \ref{thm:cdim-arithmetic-singular-ring-one-parameter-subgroups} 的记号。
对任意 \(\omega=(\omega_p)_{p\in S}\in\RR^S\)，记 \(\tau_{\omega}:\RR\to\mathfrak{R}_S\) 为与 \(\omega\) 对应的一参数子群。
取任意特征 \(\gamma\in\Gamma_S=\widehat{\mathfrak{R}_S}\)，并选取代表元 \((q,(n_p)_{p\in S})\in\ZZ[S^{-1}]\oplus\ZZ^{S}\)。
则对一切 \(t\in\RR\) 有
\begin{equation}\label{eq:cdim-arithmetic-singular-ring-character-eval}
\boxed{\;
\gamma(\tau_{\omega}(t))
=
\exp\!\Bigl(\mathrm{i}t\Bigl(q\sum_{p\in S}\omega_p+\sum_{p\in S}n_p\omega_p\Bigr)\Bigr)\; }.
\end{equation}
其中指数中的实数记为
\[
\ell_{\omega}(\gamma):=q\sum_{p\in S}\omega_p+\sum_{p\in S}n_p\omega_p\in\RR,
\]
并且该值与代表元 \((q,(n_p))\) 的选取无关。
特别地，\(\ZZ[S^{-1}]\) 分量的系数总以 \(\sum_{p\in S}\omega_p\) 的形式出现，体现了纤维积关系 \((1,-\mathbf 1_S)=0\) 在对偶端的加性约束。
\end{proposition}
\begin{proof}
由定理 \ref{thm:cdim-arithmetic-singular-ring-dual-connected}，\(\Gamma_S\cong (\ZZ[S^{-1}]\oplus\ZZ^{S})/\langle(1,-\mathbf{1}_S)\rangle\)。
对给定 \(\omega\)，令 \(\widetilde\ell_{\omega}:\ZZ[S^{-1}]\oplus\ZZ^{S}\to\RR\) 定义为
\[
\widetilde\ell_{\omega}\bigl(q,(n_p)\bigr):=q\sum_{p\in S}\omega_p+\sum_{p\in S}n_p\omega_p.
\]
则 \(\widetilde\ell_{\omega}(1,-\mathbf 1_S)=0\)，故其因子化为 \(\ell_{\omega}:\Gamma_S\to\RR\)。
Pontryagin 对偶同构 \(\Hom_{\mathrm{cont}}(\RR,\mathfrak{R}_S)\cong\Hom(\Gamma_S,\RR)\) 表明 \(\tau_{\omega}\) 对应于 \(\ell_{\omega}\)；
因此对任意 \(\gamma\in\Gamma_S\) 有 \(\gamma(\tau_{\omega}(t))=\exp(\mathrm{i}t\,\ell_{\omega}(\gamma))\)，并由代表元展开得到 \eqref{eq:cdim-arithmetic-singular-ring-character-eval}。证毕。
\end{proof}

\begin{corollary}[有限可观测投影的闭合/致密分岔]\label{cor:cdim-arithmetic-singular-ring-observable-closure-dichotomy}
取 \(\omega\in\RR^{S}\) 与有限组特征 \(\gamma_1,\dots,\gamma_m\in\Gamma_S\)，并记 \(\tau_{\omega}\) 如命题 \ref{prop:cdim-arithmetic-singular-ring-prime-frequency-character-eval}。
定义余弦读出
\[
\mathcal{O}(t):=\bigl(\Re\,\gamma_1(\tau_{\omega}(t)),\dots,\Re\,\gamma_m(\tau_{\omega}(t))\bigr)\in[-1,1]^m.
\]
令
\[
\alpha_j:=\ell_{\omega}(\gamma_j)\in\RR\qquad(1\le j\le m),
\qquad
r:=\dim_{\QQ}\mathrm{span}_{\QQ}\{\alpha_1,\dots,\alpha_m\}.
\]
则 \(\mathcal{O}(\RR)\) 的闭包等于某个 \(r\)-维子环面 \(\TT^{r}\subseteq\TT^{m}\) 在坐标实部投影下的像。
特别地：
\begin{enumerate}
  \item 若存在 \(\alpha\neq 0\) 使 \(\alpha_j\in\alpha\QQ\)（所有 \(j\)），则轨道为周期闭合，从而 \(\mathcal{O}(\RR)\) 为有限闭曲线的像；
  \item 若存在 \(j\neq k\) 使 \(\alpha_j/\alpha_k\notin\QQ\)，则相应二维投影在 \(\TT^2\) 上致密，其余弦投影给出填充型图像。
\end{enumerate}
\end{corollary}
\begin{proof}
由命题 \ref{prop:cdim-arithmetic-singular-ring-prime-frequency-character-eval}，映射
\[
t\longmapsto \bigl(\gamma_1(\tau_{\omega}(t)),\dots,\gamma_m(\tau_{\omega}(t))\bigr)
=\bigl(e^{\mathrm{i}t\alpha_1},\dots,e^{\mathrm{i}t\alpha_m}\bigr)
\]
为 \(\RR\to\TT^{m}\) 的连续群同态。
其像的闭包是由频率向量 \((\alpha_1,\dots,\alpha_m)\) 生成的最小闭子群，因而同构于维数 \(r\) 的子环面；该 \(r\) 等于 \(\{\alpha_j\}\) 的 \(\QQ\)-线性秩。
对该闭包逐坐标取实部即得 \(\mathcal{O}(\RR)\) 的闭包结论。
当全部频率共线于一维 \(\QQ\)-子空间时轨道周期闭合；若存在无理比，则二维子投影为 Kronecker 流并致密。证毕。
\end{proof}

\begin{proposition}[纤维平移的相位量子化]\label{prop:cdim-arithmetic-singular-ring-fiber-phase-quantization}
设 \(\mathcal{K}_S:=\ker(\pi_S:\Sigma_S\to\TT)\)（可同一化为 \(\prod_{p\in S}\ZZ_p\)）并将其嵌入 \(\mathfrak{R}_S\) 为 \(k\mapsto (k,(1)_{p\in S})\)。
取 \(\omega\in\RR^{S}\)、\(\tau_{\omega}:\RR\to\mathfrak{R}_S\) 以及特征 \(\gamma\in\Gamma_S\)。
设 \(\gamma\) 的 \(\ZZ[S^{-1}]\) 分量为 \(q\in\ZZ[S^{-1}]\)，并令 \(N\ge 1\) 为满足 \(Nq\in\ZZ\) 的最小整数。
则对任意 \(k\in \mathcal{K}_S\) 有 \(\gamma(k)\in\mu_N\)（\(N\)-次单位根），并且
\[
\gamma\bigl(k\cdot\tau_{\omega}(t)\bigr)=\gamma(k)\,\gamma(\tau_{\omega}(t))\qquad(\forall t\in\RR).
\]
因此存在 \(\delta_{\gamma}(k)\in\frac{2\pi}{N}\ZZ\) 使
\[
\Re\,\gamma\bigl(k\cdot\tau_{\omega}(t)\bigr)
=\cos\!\bigl(t\,\ell_{\omega}(\gamma)+\delta_{\gamma}(k)\bigr).
\]
\leanverified{paper\_cdim\_fiber\_phase\_quantization}
\end{proposition}
\begin{proof}
对任意 \(k\in\mathcal{K}_S\) 有 \(\pi_S(k)=1\)，故整数 \(n\in\ZZ\) 作为 \(\ZZ[S^{-1}]\) 中的元素在 \(\mathcal{K}_S\) 上诱导平凡特征。
因此 \(\mathcal{K}_S\) 上的特征仅依赖于 \(q\) 在商群 \(\ZZ[S^{-1}]/\ZZ\) 中的类。
由于 \(Nq\in\ZZ\)，其限制满足 \(\gamma(k)^N=\gamma(Nq,k)=1\)，从而 \(\gamma(k)\in\mu_N\)。
乘法同态性给出 \(\gamma(k\cdot\tau_{\omega}(t))=\gamma(k)\gamma(\tau_{\omega}(t))\)。
写 \(\gamma(k)=e^{\mathrm{i}\delta_{\gamma}(k)}\)（\(\delta_{\gamma}(k)\in\frac{2\pi}{N}\ZZ\)）并用 \(\gamma(\tau_{\omega}(t))=e^{\mathrm{i}t\ell_{\omega}(\gamma)}\)，取实部即得结论。证毕。
\end{proof}

\endinput
